\section*{الفصل \ref{ch:b1:complexation} --- التعقيد}

\begin{solution}{exo:b1:complexation:1}
$[\ce{Cu(NH3)4}]^{2+}$: \ce{Cu^2+}، وأربع جزيئات \ce{NH3} مرتبطة عبر N.
$[\ce{Fe(CN)6}]^{4-}$: \ce{Fe^2+}، وستة أيونات \ce{CN-} مرتبطة عبر C؛
$[\ce{Fe(CN)6}]^{3-}$: الشيء نفسه مع \ce{Fe^3+}. $[\ce{Ag(NH3)2}]^{+}$:
\ce{Ag+}، وجزيئا \ce{NH3} عبر N. $[\ce{FeSCN}]^{2+}$: \ce{Fe^3+}، وأيون
\ce{SCN-} واحد، مرتبط عبر N رغم طريقة كتابة الصيغة.
$[\ce{Zn(OH)4}]^{2-}$: \ce{Zn^2+}، وأربعة أيونات \ce{OH-} عبر O.
$[\ce{CaY}]^{2-}$: \ce{Ca^2+}، وأيون إديتا واحد عبر ذرتي N وأربع ذرات O.
\end{solution}

\begin{solution}{exo:b1:complexation:2}
$\log K_i = \log\beta_i - \log\beta_{i-1}$: 4.10، 3.41، 2.82، 2.03،
و$\mathrm{p}K_{d,i} = \log K_i$ تأخذ القيم نفسها. والأخيرة تشير إلى
$[\ce{Cu(NH3)4}]^{2+} \rightleftharpoons [\ce{Cu(NH3)3}]^{2+} + \ce{NH3}$.
\end{solution}

\begin{solution}{exo:b1:complexation:3}
المناطق: $[\ce{Cu(NH3)4}]^{2+}$ تحت pNH3 2.03، ثم $n = 3$ حتى 2.82،
و$n = 2$ حتى 3.41، و$n = 1$ حتى 4.10، و\ce{Cu^2+} فوق ذلك. pNH3 = 0:
$[\ce{Cu(NH3)4}]^{2+}$؛ 2.5: $[\ce{Cu(NH3)3}]^{2+}$؛ 3.5:
$[\ce{CuNH3}]^{2+}$؛ 5.0: \ce{Cu^2+}.
\end{solution}

\begin{solution}{exo:b1:complexation:4}
$[\ce{FeSCN^2+}]/[\ce{Fe^3+}] = \beta_1[\ce{SCN-}] = 10^{2.96} \times 0.10 =
91$: فالكسر المعقَّد $91/92 = \qty{98.9}{\percent}$.
\end{solution}

\begin{solution}{exo:b1:complexation:5}
كل $[\mathrm{ML}_k] = \beta_k[\mathrm{M}][\mathrm{L}]^k$؛ وبالقسمة على
المجموع على $k$ (الفلز الكلي) يُحذف $[\mathrm{M}]$. ومع ثلاثة
أنواع، $f_i = 1/(1 + [\mathrm{ML}_{i-1}]/[\mathrm{ML}_i] +
[\mathrm{ML}_{i+1}]/[\mathrm{ML}_i]) = 1/(1 + 1/(K_i x) + K_{i+1}x)$ مع
$x = [\mathrm{L}]$. ويكون المقام أصغر ما يمكن حين تنعدم مشتقته
$-1/(K_i x^2) + K_{i+1}$، أي $x^2 = 1/(K_iK_{i+1})$، حيث يتساوى
الجاران، و$\mathrm{pL} = \frac12(\log K_i + \log K_{i+1})$.
ومن أجل $[\ce{Cu(NH3)2}]^{2+}$: $\frac12(3.41 + 2.82) = 3.12$.
\end{solution}

\begin{solution}{exo:b1:complexation:6}
$[\ce{FeSCN}]^{2+} + \ce{F-} \rightleftharpoons [\ce{FeF}]^{2+} + \ce{SCN-}$،
$K = 10^{6.85 - 2.96} = 10^{3.89}$. والنسبة $= K[\ce{F-}]/[\ce{SCN-}] =
10^{3.89} \times 0.010/0.10 = 780$: فيُستبدل \omterm{def:b1:complexation:complex}{المعقد} الأحمر كله تقريبًا
ويبهت اللون.
\end{solution}

\begin{solution}{exo:b1:complexation:7}
$[\ce{MgY}]^{2-} + \ce{Ca^2+} \rightleftharpoons [\ce{CaY}]^{2-} + \ce{Mg^2+}$،
$K = 10^{12.69 - 10.90} = 10^{1.79} = 62$. ومع كميتين متساويتين،
$x^2/(1 - x)^2 = 62$، $x/(1 - x) = 7.9$، $x = 0.89$: فيتحرر \qty{89}{\percent} من
المغنيسيوم.
\end{solution}

\begin{solution}{exo:b1:complexation:8}
\ce{Ag2O(s) + H2O + 4NH3 <=> 2[Ag(NH3)2]+ + 2OH-}، $K = \beta_2^2 \times
10^{-15.43} = 10^{14.44 - 15.43} = 10^{-0.99}$. والثابت ليس
صغيرًا: فما إن تبلغ الأمونيا الحرة بضعة أعشار المول في
اللتر حتى يذوب الأكسيد البني المتكوّن بالقطرات الأولى (إذ تسلك الأمونيا
سلوك قاعدة) في أمين الفضة.
\end{solution}

\begin{solution}{exo:b1:complexation:9}
عند pH 10: $\alpha = 1 + 10^{1.24} + 10^{-1.96} = 18.4$، $\log\beta' = 12.69 -
1.26 = 11.43$. وعند pH 7: $\alpha = 1 + 10^{4.24} + 10^{4.04} + 10^{0.19} +
\cdots = 10^{4.45}$، $\log\beta' = 8.24$. فعند pH 7 يكون الثابت المشروط
دون $10^{10}$: فتفاعل الكالسيوم مع الإديتا ليس تامًا
بما يكفي لمعايرة، لأن أكثر الإديتا مبرتن.
\end{solution}

\begin{solution}{exo:b1:complexation:10}
1.~\ce{CuO(s) + H2O + 4NH3 <=> [Cu(NH3)4]^2+ + 2OH-}، $K = \beta_4 \times
10^{-20.64} = 10^{12.36 - 20.64} = 10^{-8.28}$.
2.~$\mathrm{pH} = 7 + \frac12(9.25 + \log 1.0) = 11.63$، $[\ce{OH-}] =
10^{-2.37}$، و$[\ce{Cu(NH3)4^2+}] = 10^{-8.28}/10^{-4.75} =
\qty{3.0e-4}{mol/L}$: لا يذوب إلا القليل جدًا.
3.~$\mathrm{pH} = 9.25$، $[\ce{OH-}] = 10^{-4.75}$، وتعطي الصيغة
$10^{-8.28 + 9.50} = 10^{1.22}$، أكثر بكثير مما تستطيع الأمونيا حمله: فيذوب
الأكسيد حتى تنفد الأمونيا، لا حتى يُبلغ التوازن. وأيونات
الأمونيوم تستهلك الهيدروكسيد الذي يحرره الذوبان؛ ولهذا
تساعد أملاح الأمونيوم الأمونيا على إذابة مركبات النحاس.
\end{solution}

\begin{solution}{exo:b1:complexation:11}
1.~التفاعل هو التكوين الثاني مطروحًا منه الأول: $K = K_2/K_1 =
10^{3.90 - 3.32} = 10^{0.58} = 3.8 > 1$، فيتفاعل $[\ce{AgNH3}]^{+}$ مع
نفسه حين يكون النوع الرئيسي.
2.~حسب \cref{exo:b1:complexation:5}، يقع أكبر كسر عند
$\mathrm{pNH_3} = \frac12(3.32 + 3.90) = 3.61$، حيث $\beta_1 x = 10^{-0.29}
= 0.51$ و$\beta_2x^2 = 1.0$: $f = 0.51/(1 + 0.51 + 1.0) = 0.20$. فلا يزيد
ما يوجد من الفضة على هيئة $[\ce{AgNH3}]^{+}$ أبدًا على \qty{20}{\percent}.
\end{solution}

\begin{solution}{exo:b1:complexation:12}
1.~\ce{CaCO3(s) + Y^4- <=> [CaY]^2- + CO3^2-}، $K = \beta K_s = 10^{12.69 -
8.30} = 10^{4.39}$.
2.~$x^2/(0.10 - x) = 10^{4.39}$ يعطي $x = \qty{0.10}{mol/L}$ في حدود
$4 \times 10^{-7}$: فتُستهلك الإديتا كلها، أي \qty{10.0}{g} من كربونات
الكالسيوم في اللتر.
3.~في الحمض تتبرتن الإديتا ($\mathrm{p}K_{a4} = 11.24$) ويهبط
ثابتها المشروط؛ وتتبرتن الكربونات أيضًا، وهذا يساعد،
لكن \omterm{def:b1:complexation:complex}{المعقد} هو ما يمسك الكالسيوم. وإبقاء pH مرتفعًا يُبقي
الإديتا على هيئة \ce{Y^4-}.
\end{solution}

\begin{solution}{pb:b1:complexation:1}
\textbf{1.} \ce{Ag+ + NH3 <=> [AgNH3]+}، $\beta_1 = [\ce{AgNH3+}]/([\ce{Ag+}][\ce{NH3}])$؛
\ce{Ag+ + 2NH3 <=> [Ag(NH3)2]+}، $\beta_2 = [\ce{Ag(NH3)2+}]/([\ce{Ag+}][\ce{NH3}]^2)$.
\textbf{2.} $\log K_1 = 3.32$، $\log K_2 = 7.22 - 3.32 = 3.90$.
\textbf{3.} كان $[\ce{AgNH3}]^{+}$ سيغلب بين pNH3 3.90 و3.32،
وهي فترة معكوسة الاتجاه: فليست له منطقة.
\textbf{4.} \ce{2[AgNH3]+ <=> Ag+ + [Ag(NH3)2]+}، $K = K_2/K_1 = 10^{0.58} =
3.8$.
\textbf{5.} $[\ce{Ag(NH3)2}]^{+}$ تحت $\frac12\log\beta_2 = 3.61$ و\ce{Ag+}
فوقها.
\textbf{6.} $\beta_2[\ce{NH3}]^2 = 10^{7.22 - 2.00} = 10^{5.22} = \num{1.7e5}$.
\textbf{7.} $\mathrm{p}K_d = \log\beta_2 = 7.22$.
\textbf{8.} \ce{AgCl(s) + 2NH3 <=> [Ag(NH3)2]+ + Cl-}، $K = \beta_2K_s =
10^{7.22 - 9.75} = 10^{-2.53}$.
\textbf{9.} $s = \sqrt{K}\,[\ce{NH3}] = 10^{-1.265} \times 1.0 =
\qty{0.054}{mol/L}$.
\textbf{10.} في الماء $s = 10^{-4.875} = \qty{1.3e-5}{mol/L}$: أي أكثر بنحو 4000
مرة.
\textbf{11.} $0.0543 \times 143.4 = \qty{7.8}{g}$.
\textbf{12.} $[\ce{Ag+}] = K_s/[\ce{Cl-}] = 10^{-9.75}/0.054 =
\qty{3.3e-9}{mol/L} \ll s$.
\textbf{13.} $s = \sqrt{K}(1.0 - 2s)$، $s = 0.0543/(1 + 2 \times 0.0543) =
\qty{0.049}{mol/L}$.
\textbf{14.} pH الأمونيا ذات التركيز \qty{1.0}{mol/L} يساوي 11.6، وهو أكثر من $9.25 + 1$:
فلا يكون منها على هيئة \ce{NH4+} إلا \qty{0.4}{\percent}.
\textbf{15.} $K = 10^{7.22 - 12.27} = 10^{-5.05}$.
\textbf{16.} $s = 10^{-2.525} = \qty{3.0e-3}{mol/L}$، أي \qty{0.56}{g/L}.
\textbf{17.} $K = 10^{-8.85}$، $s = 10^{-4.425} = \qty{3.8e-5}{mol/L}$،
أي \qty{8.8}{mg/L}.
\textbf{18.} عالج كل راسب بأمونيا مخففة: لا يذوب إلا الكلوريد. وتذيب
الأمونيا المركزة البروميد، ببطء أكبر وجزئيًا؛ أما اليوديد فلا يذوب.
\textbf{19.} $s = 1.0/187.8 = \qty{5.3e-3}{mol/L}$؛ $[\ce{NH3}] = s/\sqrt{K}
= \num{5.3e-3}/10^{-2.525} = \qty{1.8}{mol/L}$.
\textbf{20.} $s = \qty{4.3e-3}{mol/L}$ و$[\ce{NH3}] = \num{4.3e-3}/10^{-4.425}
= \qty{113}{mol/L}$، أي ضعف تركيز الماء في الماء النقي
(\qty{55}{mol/L}): وهذا مستحيل. فالأمونيا لا تستطيع إذابة يوديد الفضة.
\textbf{21.} \ce{[Ag(NH3)2]+ + Cl- + 2H3O+ -> AgCl(s) + 2NH4+ + 2H2O}،
$K = 10^{2.53} \times (10^{9.25})^2 = 10^{21.03}$: يعود الراسب الأبيض
إلى الظهور، وهذا يؤكد أنه كلوريد الفضة.
\textbf{22.} $s = 1.0/143.4 = \qty{6.97e-3}{mol/L}$؛ $[\ce{NH3}] =
s/\sqrt{K} = \num{6.97e-3}/10^{-1.265} = \qty{0.128}{mol/L}$.
\textbf{23.} $2s = \qty{0.0139}{mol/L}$.
\textbf{24.} $[\ce{Ag+}] = 10^{-9.75}/\num{6.97e-3} = \qty{2.6e-8}{mol/L}$،
وهو مهمل بجانب $s$.
\textbf{25.} $0.128 + 0.014 = $ \textbf{\qty{0.142}{mol/L} من الأمونيا}:
فالكمية \qty{0.142}{mol} في اللتر تذيب \qty{1.0}{g} من كلوريد
الفضة بالضبط.
\end{solution}
